\documentclass[11pt,a4paper]{article}

% ---------------------------------------------------------------
% Packages (same set-up as the THM model report)
% ---------------------------------------------------------------
\usepackage[T1]{fontenc}
\usepackage[utf8]{inputenc}
\usepackage{lmodern}
\usepackage[a4paper,margin=2.5cm]{geometry}
\usepackage{amsmath,amssymb,bm}
\usepackage{booktabs,tabularx,array}
\usepackage{siunitx}
\usepackage{enumitem}
\usepackage{xcolor}
\usepackage{graphicx}
\usepackage{multirow}
\usepackage[numbers,sort&compress]{natbib}
\usepackage[colorlinks=true,linkcolor=blue!50!black,citecolor=blue!50!black,urlcolor=blue!50!black]{hyperref}

\sisetup{range-phrase=--, range-units=single}
\newcolumntype{L}{>{\raggedright\arraybackslash}X}
\renewcommand{\arraystretch}{1.2}
\graphicspath{{../figures/}}
\newcommand{\open}{$\square$}
\newcommand{\incons}{\textcolor{red!70!black}{$\ddagger$}}
\newcommand{\added}{\textcolor{blue!60!black}{$^{+}$}}

% ---------------------------------------------------------------
\title{GREAT Cell Modelling Specification\\[0.5ex]
\large DECOVALEX 2027 -- SAFENET-2}
\author{Peng-Zhi Pan$^{1,2}$, Wenbo Hou$^{1,2}$, Ning Yang$^{1,2}$,
Mostafa Mollaali$^{3,4}$, Olaf Kolditz$^{4,5}$\\[2ex]
\small $^{1}$State Key Laboratory of Geomechanics and Geotechnical Engineering Safety,\\
\small Institute of Rock and Soil Mechanics, Chinese Academy of Sciences (CAS), Wuhan, China\\
\small $^{2}$University of Chinese Academy of Sciences, Beijing, China\\
\small $^{3}$Technische Universit\"at Bergakademie Freiberg (TUBAF), Germany\\
\small $^{4}$Helmholtz Centre for Environmental Research -- UFZ, Leipzig, Germany\\
\small $^{5}$Technische Universit\"at Dresden, Germany}
\date{Version 2 -- 22 November 2026}

\begin{document}
\maketitle

\begin{abstract}
This specification fixes geometry, material parameters, loads, boundary
conditions and outputs for modelling the GREAT cell experiments, so that all
DECOVALEX SAFENET-2 teams simulate the same problem and their results can be
compared directly. It is code-independent. Version~2 defines two thermal
scenarios for two application areas: Scenario~I, heating of the sample from
its boundary followed by a short cold-water injection (the scenario of
version~1), and Scenario~II, a geothermal application with cold-water
injection into a hot (\SI{300}{\celsius}) rock. OpenGeoSys (OGS) reference
solutions exist for the cases M1 and THM2 of Scenario~I.
\end{abstract}

\tableofcontents

% ---------------------------------------------------------------
\section{Purpose and scope}
% ---------------------------------------------------------------
\subsection{General rules}
\begin{itemize}
  \item \textbf{Fixed:} everything in Sections~\ref{sec:geometry} to
        \ref{sec:outputs}. Deviations must be reported.
  \item \textbf{Free:} numerical method, element type, mesh, time stepping and
        solver, within the accuracy targets of Section~\ref{sec:verification}.
  \item \textbf{Units:} SI throughout (m, s, Pa, K, kg). Strains as fractions;
        plots may use $10^{-2}$\,\%. Exception: the Scenario~II outputs use
        \si{\celsius}, MPa (gauge), ml/min and ml (Section~\ref{sec:outputs-II}).
  \item \textbf{Model dimension:} 2D plane strain, horizontal mid-section of
        the sample, unit thickness \SI{1}{m}. All rates are per metre thickness;
        in Scenario~II the 3D injection rate is divided by the specimen thickness
        $B=\SI{0.2}{m}$ (Section~\ref{sec:schedule-II}).
\end{itemize}

\subsection{Two thermal scenarios}
\label{sec:scenarios}
The thermal benchmarks are run in two scenarios that represent two application
areas (Table~\ref{tab:scenarios}). Both use the same disc geometry, angle convention, mechanical load
cases and rock mechanical parameters; they differ in the rubber sheath, in how
heat enters the sample, in the temperature range, the injection programme and
the hydraulic parameters.

\begin{itemize}
  \item \textbf{Scenario~I -- heating of the sample (repository-type
        application).} The sample starts at laboratory temperature and is heated
        from its outer boundary, as by the heating bands of the GREAT cell, while
        it is held under the poly-axial load. A short cold-water injection at the
        centre follows. Heat enters by conduction from the boundary; the
        scenario tests thermal expansion under poly-axial stress, as in the near
        field of a heat-generating waste repository. This is the scenario of
        version~1 of this specification; it is kept unchanged.
  \item \textbf{Scenario~II -- cold-water injection into a hot rock (geothermal
        application).} Rock and pore water start uniformly at
        \SI{300}{\celsius}. After the poly-axial load is applied, water at
        \SI{15}{\celsius} is injected at a constant low rate for \SI{3600}{s},
        while the outer boundary is thermally insulated. The injected water is
        the only thermal input, so the scenario tests cooling-induced
        contraction, advective heat transport and pressure propagation, as
        around an injection well in an enhanced geothermal system. It follows
        the SAFENET-2 THM benchmark proposal of~\cite{Pan2026}.
\end{itemize}

\begin{table}[htbp]
\centering
\caption{Overview of the two thermal scenarios.}
\label{tab:scenarios}
\small
\begin{tabularx}{\textwidth}{@{}lLL@{}}
\toprule
Item & Scenario I: heating (version 1) & Scenario II: geothermal \\
\midrule
Application        & Repository near field (heat-generating waste) & Enhanced geothermal system, injection well \\
Rubber sheath      & Yes: \SI{3}{mm}, $0.097 < r \le \SI{0.100}{m}$ & No: rock up to $R=\SI{0.100}{m}$ \\
Initial temperature & \SI{20}{\celsius} (\SI{293}{K}) & \SI{300}{\celsius} uniform, $T_\mathrm{ref}=T_0$ \\
Outer thermal boundary & Heated: \SI{20}{\celsius} $\rightarrow$ \SI{90}{\celsius} in \SI{325}{s}, then held & Adiabatic \\
Injected water     & \SI{20}{\celsius}, \SI{1}{l/min} per m, 4000--4050\,s & \SI{15}{\celsius}, $Q_0=\SI{0.1}{ml/min}$, ramp 2000--2120\,s, constant to 5600\,s \\
Mechanical load    & Load~A applied at $t=0$ & Load ramped 0--1800\,s, then held \\
Pore pressure      & $p_0=\SI{1}{MPa}$; outer boundary $p=p_0$ & Gauge pressure $p_0=0$; outer boundary $p=0$ (absolute reference \SI{8.6}{MPa}) \\
Matrix permeability & \SI{1.409e-14}{m^2} & \SI{2.58e-19}{m^2} (Greywacke), \SI{1e-19}{m^2} (Gneiss) \\
Water properties   & Constant (compressibility \SI{4.4e-10}{Pa^{-1}}) & Temperature-dependent (IAPWS-IF97) \\
Duration           & \SI{5000}{s} & \SI{5600}{s} \\
Isothermal control & -- & HM run at \SI{15}{\celsius}, same mesh, time steps and flow schedule \\
Case ID (example)  & THM2, Load A, Greywacke (version 1 naming) & THM2-GW-A-T300~\cite{Pan2026} \\
\bottomrule
\end{tabularx}
\end{table}

Sections~\ref{sec:geometry} to \ref{sec:outputs} cover both scenarios; where
they differ, subsections or table rows are marked Scenario~I or~II. The
verification values of Section~\ref{sec:verification} are so far available for
Scenario~I only.

% ---------------------------------------------------------------
\section{Benchmark matrix}
% ---------------------------------------------------------------
The benchmarks add one process at a time: mechanics (M), then fluid flow (HM),
then heat (THM), each for an intact sample, a single fracture and a fracture
network (Fig.~\ref{fig:concept}, Table~\ref{tab:matrix}). For Scenario~I this
version specifies \textbf{M1} and \textbf{THM2} in full; the fractured cases
follow the same conventions and are added later.

Scenario~II adopts the THM case matrix of~\cite{Pan2026} (Table~5 there),
restricted to the initial temperature $T_0=\SI{300}{\celsius}$. Each THM run is
paired with an isothermal HM control at \SI{15}{\celsius} with the same mesh,
time steps and flow schedule; the controls are not included in the run counts.
The runs are carried out in two phases: Phase~1 contains the mandatory runs for
Greywacke under Load~A; the mandatory suite adds Loads B and C for the
fractured samples. THM4, THM5, Loads E and F, and Gneiss are optional. Case IDs
follow the format \texttt{THM\{number\}-\{rock\}-\{load\}-T\{$T_0$\}}, e.g.\
\texttt{THM3b-GW-A-T300}, with GW for Greywacke and GN for Gneiss; the HM
control of this run is \texttt{HM3b-GW-A-T15}.

\begin{figure}[htbp]
\centering
\includegraphics[width=0.6\textwidth]{great-benchmarking-thm.png}
\caption{GREAT cell benchmark concept: M (grey), HM (blue, fluid injection) and
THM (red boundary, heating) for intact, fractured and networked samples.}
\label{fig:concept}
\end{figure}

\begin{table}[htbp]
\centering
\caption{Benchmark cases. Run counts for Scenario~II: Phase~1 / mandatory suite
/ full suite, for $T_0=\SI{300}{\celsius}$, without the HM controls.}
\label{tab:matrix}
\small
\begin{tabularx}{\textwidth}{@{}lLLcL@{}}
\toprule
Case & Sample / configuration & Processes, loads & Runs & Status \\
\midrule
\multicolumn{5}{@{}l}{\textit{Scenario I -- heating (version 1)}} \\
M1             & Intact                    & Poly-axial load, rotated (Loads A, E, B, F, C) & -- & Specified; OGS reference \\
M3a / M3b      & Full / half fracture      & Mechanics                                      & -- & Geometry given; load as M1 \\
HM2            & Intact, central borehole  & Load + injection                               & -- & Covered by THM2 without heating \\
HM3a / HM3b    & Full / half fracture      & Load + injection into fracture                 & -- & To be specified \\
HM4            & Fracture network          & Load + injection                               & -- & To be specified \\
THM2           & Intact                    & Load A + heating + cold-water injection        & -- & Specified; OGS reference \\
THM3a/b, THM4  & Fractured                 & As HM3/HM4 + heating                           & -- & To be specified \\
\midrule
\multicolumn{5}{@{}l}{\textit{Scenario II -- geothermal, $T_0=\SI{300}{\celsius}$, Greywacke and Gneiss~\cite{Pan2026}}} \\
THM2           & Intact disk; centre-point injection                     & Load + cold-water injection; A               & 1 / 1 / 2  & Mandatory (GW, A) \\
THM3a          & Through-going fracture, length \SI{0.188}{m}            & As THM2, injection into fracture; A/E/B/F/C  & 1 / 3 / 10 & Mandatory (GW; A, B, C) \\
THM3b          & Embedded (half) fracture, length \SI{0.04}{m}           & As THM3a; A/E/B/F/C                          & 1 / 3 / 10 & Mandatory (GW; A, B, C) \\
THM4           & Propagating fracture; initial crack \SI{0.02}{m} through the centre, inclined at \ang{45} (HM3c) & As THM2, injection into the crack centre; A/B/C & 0 / 0 / 6  & Optional \\
THM5           & Centre borehole, radius \SI{0.0025}{m}, no notch (HM3d) & As THM2, fracture nucleation at the borehole; A/B/C & 0 / 0 / 6 & Optional \\
\midrule
Total (Sc.~II) &                                                         &                                              & 3 / 7 / 34 & Each with an HM control at \SI{15}{\celsius} \\
\bottomrule
\end{tabularx}
\end{table}

In Scenario~II the names follow~\cite{Pan2026}: THM4 denotes a propagating
fracture and THM5 fracture initiation from a centre borehole, whereas in
Scenario~I and Fig.~\ref{fig:concept} HM4/THM4 denote a fracture network.
The Scenario~II cases correspond to the HM benchmarks of~\cite{Mollaali2026}
as follows: THM2 -- HM1, THM3a -- HM2a, THM3b -- HM2b, THM4 -- HM3c,
THM5 -- HM3d (``HM5'' in~\cite{Pan2026}).
\cite{Pan2026} also defines the initial temperature \SI{90}{\celsius}, which
doubles its run counts (6 / 14 / 68); it is not part of this specification.

Teams are encouraged to run the stepwise sequence T $\rightarrow$ TH
$\rightarrow$ THM used for the reference model, since each step has its own
check.

% ---------------------------------------------------------------
\section{Geometry and conventions}
\label{sec:geometry}
% ---------------------------------------------------------------
The domain is a disc of radius $R=\SI{0.100}{m}$ centred at $(0,0)$
(Fig.~\ref{fig:geometry}, Table~\ref{tab:geometry}). In Scenario~I the rock
fills $r\le\SI{0.097}{m}$ and a \SI{3}{mm} rubber sheath lies outside it, as
in the GREAT cell. In Scenario~II there is no sheath: the rock fills the whole
disc and the PEE/DSS loads act directly on the rock surface. All other
geometric items are the same in both scenarios. The sample radius
$R=\SI{0.100}{m}$ and the through-going fracture length of \SI{0.188}{m} apply
to both scenarios; the values \SI{193.75}{mm} (sample diameter in Fig.~3)
and \SI{0.194}{m} (THM3a fracture length) of~\cite{Pan2026} are not used. Angles $\theta$ are measured
counter-clockwise from the $+x$ axis (Section~\ref{sec:angles}).

\begin{figure}[htbp]
\centering
\includegraphics[width=\textwidth]{fig_geometry.pdf}
\caption{Geometry and PEE numbering (a), detail of the sheath at PEE~1a (b), shown for Scenario~I. In Scenario~II the sheath is replaced by rock.}
\label{fig:geometry}
\end{figure}

\begin{table}[htbp]
\centering
\caption{Geometry and conventions.}
\label{tab:geometry}
\begin{tabularx}{\textwidth}{@{}lL@{}}
\toprule
Item & Specification \\
\midrule
Rock               & Scenario~I: $r \le \SI{0.097}{m}$ (material 0). Scenario~II: $r \le \SI{0.100}{m}$ (material 0) \\
Sheath             & Scenario~I only: $\SI{0.097}{m} < r \le \SI{0.100}{m}$ (material 1) \\
Centre point       & $(0,0)$; injection point (THM2), borehole not resolved \\
Probe circle       & $r=\SI{0.065}{m}$; volumetric strain output \\
PEE arcs           & 16 arcs of \ang{18} on $r=\SI{0.100}{m}$. PEE $i$ centred at $\theta=-(i-1)\cdot\ang{22.5}$ ($i=1\ldots8$, clockwise from $+x$); PEE $i$a opposite, at $+\ang{180}$ \\
DSS arcs           & 16 arcs of \ang{4.5} between the PEEs. DSS $i$ lies between PEE $i$ and PEE $i+1$; DSS~8 between PEE~8 and PEE~1a \\
Rigid-body control & $P_\mathrm{bottom}=(0,\,-0.09894)$\,m: $u_x=u_y=0$. $P_\mathrm{left}=(-0.09894,\,0)$\,m: $u_y=0$ \\
Full fracture (M3a/HM3a/THM3a) & Segment $(-0.094,\,0.094)\,\mathrm{m}\times\{0\}$ \\
Half fracture (M3b/HM3b/THM3b) & Segment $(0,\,0.04)\,\mathrm{m}\times\{0\}$ \\
Initial crack (THM4, Sc.~II)\incons & Segment through the centre, $(x,y)=s\,(\cos\ang{45},\sin\ang{45})$, $s\in[-0.01,\,0.01]$\,m, i.e.\ length \SI{0.02}{m} at $\theta=\ang{45}$ (HM3c of~\cite{Mollaali2026}); \cite{Pan2026} states \ang{-45} \\
Borehole (THM5, Sc.~II) & Centre borehole, radius \SI{0.0025}{m}, no notch (HM3d of~\cite{Mollaali2026}) \\
\bottomrule
\end{tabularx}
\end{table}

The PEE and DSS arc widths (\ang{18}/\ang{4.5}) are estimates from Fig.~4b
of~\cite{Mollaali2026}; teams should use them until measured values are
available. The rigid-body points lie in the sheath in Scenario~I and in the rock
in Scenario~II; their coordinates are the same.

\subsection{Angle convention}
\label{sec:angles}
This specification keeps its original angle convention (Fig.~\ref{fig:angles}a):
$\theta$ is measured counter-clockwise from the $+x$ axis, PEE~1 is centred at
$\theta=\ang{0}$ and the PEEs are numbered clockwise. The monitoring points
P0, P90, P180 and P270 of~\cite{Pan2026} (Table~10 there) are given in the same
convention, e.g.\ P0 $=(+0.065,\,0)$\,m.

The laboratory schematic of the GREAT cell (Fig.~3 of~\cite{Pan2026}) uses a
different convention (Fig.~\ref{fig:angles}b): the angle $\varphi$ is measured
clockwise from the top of the drawing, with PEE~2 at $\varphi=\ang{0}$ and
PEE~1 at $\varphi=\ang{337.5}$. Both conventions describe the same PEE
arrangement; they differ in the zero direction and the sense of rotation:
\begin{equation}
  \varphi = (-\theta - \ang{22.5}) \bmod \ang{360}, \qquad
  \theta = (-\varphi - \ang{22.5}) \bmod \ang{360}.
\end{equation}
Fig.~\ref{fig:angles}b is for explanation only. Profiles plotted against
$\varphi$ instead of $\theta$ appear mirrored and shifted by \ang{22.5},
which corresponds to one PEE spacing and to the difference between Load~A and
Load~E. All inputs and outputs of this specification use $\theta$.

\begin{figure}[htbp]
\centering
\includegraphics[width=\textwidth]{fig_angle_conventions.pdf}
\caption{Angle conventions. (a) Convention of this specification (used for
all inputs and outputs), with the monitoring points P0--P270 on the probe
circle. (b) Convention of the laboratory schematic in~\cite{Pan2026}, Fig.~3,
for explanation only. Both panels show the same PEE arrangement; (b) is drawn
with PEE~2 at the top, as in the laboratory schematic.}
\label{fig:angles}
\end{figure}

Note: in Fig.~3 of~\cite{Pan2026} the labels 2A and 3A appear swapped; with
diametrically opposite pairs, 2A lies at $\varphi=\ang{180}$ and 3A at
$\varphi=\ang{202.5}$, as drawn in Fig.~\ref{fig:angles}b.

% ---------------------------------------------------------------
\section{Mesh}
\label{sec:mesh}
% ---------------------------------------------------------------
The mesh is the team's choice, but it must resolve four features: the sheath
(Scenario~I), the PEE/DSS arc ends, the probe circle and the centre point.
\begin{itemize}
  \item PEE/DSS arc ends and the two fixed points must be mesh nodes.
  \item The probe circle should be embedded as mesh edges, or outputs
        interpolated onto it at $\le\ang{1}$ spacing.
  \item Scenario~I: at least two element layers across the \SI{3}{mm} sheath.
  \item Refinement towards the centre for injection; the reference uses
        $h=\SI{1.5}{mm}$ at the centre, \SI{6}{mm} at $r=\SI{0.065}{m}$ and
        \SI{2}{mm} in the sheath.
  \item Reference mesh: 11\,889 nodes, 5\,804 quadratic triangles; displacement
        quadratic, pressure and temperature linear.
\end{itemize}

% ---------------------------------------------------------------
\section{Material parameters}
\label{sec:materials}
% ---------------------------------------------------------------
All materials are linear elastic and isotropic, and mechanical properties do
not depend on temperature. Table~\ref{tab:params} lists the rock and sheath
parameters, Table~\ref{tab:water} the water parameters. Both tables compare
version~1 of this specification (Scenario~I, rock values from Table~1
of~\cite{Mollaali2026}, hydraulic and thermal values from the reference TH
model) with the THM proposal of~\cite{Pan2026} (Tables~2, 4, 6 and~7 there),
which defines Scenario~II. Where a value differs between the two sources, the
row gives both and is marked \incons; parameters added from~\cite{Pan2026} are
marked \added. Values marked ``I'' and ``II'' apply to Scenario~I and~II only.

\begin{table}[htbp]
\centering
\caption{Rock and sheath parameters. \incons: inconsistent between
version~1 and~\cite{Pan2026}; \added: added from~\cite{Pan2026}; (I), (II):
applies to Scenario~I or~II only.}
\label{tab:params}
\footnotesize
\begin{tabularx}{\textwidth}{@{}>{\raggedright\arraybackslash}p{3.3cm} l *{3}{>{\centering\arraybackslash}p{1.55cm}} L@{}}
\toprule
Parameter & Unit & Greywacke & Gneiss & Sheath (I) & Remark \\
\midrule
\multicolumn{6}{@{}l}{\textit{Mechanical}} \\
Young's modulus $E$          & GPa           & 26.85 & 83.9 & 0.1 & Consistent; Gneiss alt.\ 57.1 for sensitivity \\
Poisson's ratio $\nu$        & --            & 0.27  & 0.21 & 0.4 & Consistent; Gneiss alt.\ 0.19 \\
Solid density $\rho_s$       & kg\,m$^{-3}$  & 2650  & 2750 & 1200 & Consistent \\
Tensile strength\added       & MPa           & 16.67--17.01 & 16.8 / 6.2 & -- & \cite{Mollaali2026}: Gneiss 16.8 only; source of 6.2 to be clarified \\
Fracture energy $G_c$\added  & N\,m$^{-1}$   & 30  & 50  & -- & Phase-field models \\
Normal stiffness $k_n$\added & GPa\,m$^{-1}$ & 100 & 200 & -- & Interface and contact models \\
Shear stiffness $k_t$\added  & GPa\,m$^{-1}$ & 50  & 100 & -- & Interface and contact models \\
\midrule
\multicolumn{6}{@{}l}{\textit{Hydraulic}} \\
Porosity $\phi$              & --    & 0.005 & 0.001 & as rock & Consistent \\
Biot coefficient $\alpha_B$  & --    & 0.8   & 0.6   & 0 & Consistent \\
Permeability $k$ (I)\incons  & m$^2$ & \num{1.409e-14} & \num{1.409e-14} & as rock & From the STIMTEC-calibrated TH model; sheath in~\cite{Mollaali2026}: \num{1e-19} \\
Permeability $k$ (II)\incons & m$^2$ & \num{2.58e-19} & \num{1e-19} & -- & Gneiss range \num{1e-18} to $<$\num{1e-19}; five orders of magnitude below (I) \\
\midrule
\multicolumn{6}{@{}l}{\textit{Thermal}} \\
Bulk thermal conductivity $\lambda$ (I)\incons & W\,m$^{-1}$\,K$^{-1}$ & 2.8 & 2.8 & 2.8 & (I) prescribes bulk values \\
Solid thermal conductivity $\lambda_s$ (II)\incons & W\,m$^{-1}$\,K$^{-1}$ & 2.50 & 3.00 & -- & (II) prescribes solid values; bulk from Eq.~\eqref{eq:mixing} \\
Bulk heat capacity $\rho c$ (I)\incons & J\,m$^{-3}$\,K$^{-1}$ & \num{2.133e6} & \num{2.133e6} & \num{2.133e6} & $2700\times790$; not based on $\rho_s$ \\
Solid specific heat $c_s$ (II)\incons & J\,kg$^{-1}$\,K$^{-1}$ & 850 & 800 & -- & $\rho_s c_s=\num{2.2525e6}$ / \num{2.2e6} J\,m$^{-3}$\,K$^{-1}$ \\
Linear thermal expansion $\alpha_T$ (I)\incons & K$^{-1}$ & \num{8e-6} & \num{8e-6} & \num{8e-6} & Gneiss differs between (I) and (II) \\
Linear thermal expansion $\alpha_T$ (II)\incons & K$^{-1}$ & \num{8e-6} & \num{7e-6} & -- & \\
Reference temperature $T_\mathrm{ref}$\added & \si{\celsius} & (I) 20, (II) 300 & (I) 20, (II) 300 & 20 & $T_\mathrm{ref}=T_0$: no initial thermal stress \\
\bottomrule
\end{tabularx}
\end{table}

In Scenario~II the bulk properties follow from porosity-weighted averaging of
solid and water values~\cite{Pan2026}:
\begin{equation}
  (\rho c)_\mathrm{eff} = \phi\,(\rho c)_f + (1-\phi)\,\rho_s c_s, \qquad
  \lambda_\mathrm{eff} = \phi\,\lambda_f + (1-\phi)\,\lambda_s .
  \label{eq:mixing}
\end{equation}
In Scenario~I bulk $\lambda$ and $\rho c$ are prescribed; teams back-calculate
solid values for their own mixing rule. Storage follows in both scenarios from
$\phi$, $\alpha_B$, the water compressibility and the elasticity.

\subsection{Fracture and numerical parameters}
\label{sec:fracture}
For the fractured cases, \cite{Pan2026} takes the initial fracture state, the
aperture--permeability relation and the phase-field settings from the HM
benchmarks (``copy HM \ldots, do not refit''). Table~\ref{tab:fracture} lists
these values from~\cite{Mollaali2026}; they apply to both scenarios.

\begin{table}[htbp]
\centering
\caption{Fracture and numerical parameters of the HM benchmarks~\cite{Mollaali2026}
(Tables~1, 2 and Section~3 there). GW/GN: Greywacke/Gneiss. \incons: inconsistent
within~\cite{Mollaali2026} or with~\cite{Pan2026}.}
\label{tab:fracture}
\footnotesize
\begin{tabularx}{\textwidth}{@{}>{\raggedright\arraybackslash}p{3.6cm} l >{\raggedright\arraybackslash}p{3.4cm} L@{}}
\toprule
Parameter & Unit & Value & Remark \\
\midrule
\multicolumn{4}{@{}l}{\textit{Initial fracture state}} \\
Initial fracture width $w_\mathrm{ini}$ & m & \num{1e-6} & Static fractures (THM3a/b) and initial crack (THM4) \\
Initial pore pressure & MPa & 0.1 & HM benchmarks; Scenario~II uses gauge $p=0$ (Table~\ref{tab:icbc-II}) \\
\midrule
\multicolumn{4}{@{}l}{\textit{Fracture flow}} \\
Fracture permeability (LIE) & m$^2$ & $w^2/12$ & Cubic law; storage $S_\mathrm{fr}=\phi_f/\kappa_f$ ($\phi_f$ not given) \\
Permeability tensor (VPF) & m$^2$ & $K_m\mathbf I + (1-v)\,\xi\,(w^2/12-K_m)(\mathbf I-\mathbf n_\Gamma\otimes\mathbf n_\Gamma)$ & $\xi\ge1$ (value not given) \\
Aperture (VPF) & m & $w\approx h_e\,\varepsilon_1$ & $h_e$ element size, $\varepsilon_1$ max.\ principal strain \\
\midrule
\multicolumn{4}{@{}l}{\textit{Fracture mechanics}} \\
Fracture energy $G_c$ & N\,m$^{-1}$ & GW 30, GN 50, sheath 0.1 & Consistent with Table~\ref{tab:params} \\
Normal stiffness $k_{nn}$ (LIE) & GPa\,m$^{-1}$ & GW 100, GN 200 & \\
Shear stiffness $k_{tt}$ (LIE)\incons & GPa\,m$^{-1}$ & GW 50, GN 100 & Table~2 of~\cite{Mollaali2026} gives 100/200 for M1/HM1 \\
Strain energy split (VPF) & -- & Volumetric--deviatoric (M2a, HM3c, HM3d); isotropic (HM2a, HM2b) & AT1/AT2 choice not stated per case \\
\midrule
\multicolumn{4}{@{}l}{\textit{Discretization (VPF; LIE uses $h_\mathrm{min}=\SI{2e-3}{m}$)}} \\
Intact (THM2 / HM1) & m & $h_\mathrm{min}=\num{2e-3}$, $\ell=\num{4e-3}$ & $\ell/h_\mathrm{min}=2$ \\
Static fractures (THM3a/b / HM2a/b) & m & $h_\mathrm{min}=\num{2.5e-4}$, $\ell=\num{5e-4}$ & \\
Propagation (THM4, THM5 / HM3c, HM3d) & m & $h_\mathrm{min}=\num{5e-4}$, $\ell=\num{1e-3}$ & \\
\midrule
\multicolumn{4}{@{}l}{\textit{Injection in the HM benchmarks (for reference)}} \\
HM1 (THM2) & m$^3$\,s$^{-1}$ & \num{2.085e-9} at the centre & Outer boundary $p=0$ \\
HM2a (THM3a) & m$^3$\,s$^{-1}$ & \num{4.167e-7} (25 ml/min) at $(0,0)$ & $p=\SI{3.45}{MPa}$ at both fracture ends $(\pm0.094,0)$ \\
HM2b (THM3b) & m$^3$\,s$^{-1}$ & \num{4.167e-7} at $(0,0)$ & $p=\SI{3.45}{MPa}$ at $(0.04,0)$ \\
HM3c (THM4) & m$^3$\,s$^{-1}$ & \num{5e-8}, line element at crack centre & $q=\SI{1e-4}{m^2\,s^{-1}}$ over $h=\SI{5e-4}{m}$ \\
HM3d (THM5) & m$^3$\,s$^{-1}$ & \num{3.75e-7} around the borehole wall & $q_0=\SI{2.4e-5}{m\,s^{-1}}$ over $2\pi r$ \\
\bottomrule
\end{tabularx}
\end{table}

All HM benchmarks equilibrate under load for \SI{3000}{s} before a \SI{500}{s}
injection. In Scenario~II the schedule of Section~\ref{sec:schedule-II} applies
instead, with $Q_\mathrm{2D}=\SI{8.333e-9}{m^3\,s^{-1}}$ per m for all cases.

\begin{table}[htbp]
\centering
\caption{Water parameters. \incons: inconsistent between version~1
and~\cite{Pan2026} or within~\cite{Pan2026}; \added: added
from~\cite{Pan2026}. The last column gives IAPWS-IF97 values at
\SI{300}{\celsius} and \SI{8.6}{MPa} absolute pressure for orientation.}
\label{tab:water}
\footnotesize
\begin{tabularx}{\textwidth}{@{}>{\raggedright\arraybackslash}p{3.1cm} l >{\raggedright\arraybackslash}p{2.0cm} L c L@{}}
\toprule
Parameter & Unit & Scenario I & Scenario II~\cite{Pan2026} & IF97, \SI{300}{\celsius} & Remark \\
\midrule
Density $\rho_f$             & kg\,m$^{-3}$ & 1000 at $p_0=\SI{1}{MPa}$ & 1000; variable above \SI{100}{\celsius} & 712 & Consistent below \SI{100}{\celsius} \\
Compressibility $\beta_p$\incons & Pa$^{-1}$ & \num{4.4e-10} & 0; variable above \SI{100}{\celsius} & \num{3.2e-9} & I includes it (storage); II neglects it up to \SI{100}{\celsius} \\
Dynamic viscosity $\mu$      & Pa\,s & \num{1e-3} & \num{1e-3}; variable above \SI{100}{\celsius} & \num{8.6e-5} & Consistent below \SI{100}{\celsius} \\
Specific heat $c_f$          & J\,kg$^{-1}$\,K$^{-1}$ & 4180 & 4180; variable above \SI{100}{\celsius} & 5751 & Consistent below \SI{100}{\celsius} \\
Thermal conductivity $\lambda_f$\incons & W\,m$^{-1}$\,K$^{-1}$ & 0.6 & 0.6, constant & 0.553 & II keeps 0.6 at \SI{300}{\celsius}, 8\,\% above IF97 \\
Thermal expansion $\beta_f$\incons & K$^{-1}$ & 0 (neglected) & 0; variable above \SI{100}{\celsius} & \num{3.3e-3} & Large at \SI{300}{\celsius}; drives pressure change on cooling \\
Injection temperature $T_\mathrm{in}$\incons & \si{\celsius} & 20 & 15, only during inflow & -- & Scenario-specific \\
Property model\added         & -- & Constant & IAPWS-IF97; $\rho_f$, $c_f$ of $T$ and $p$, $\mu$ of $T$ & -- & \\
Pressure reference\added     & MPa & $p_0=1$ (absolute) & Gauge $p$; $p_\mathrm{abs}=p+8.6$ at \SI{300}{\celsius} & -- & Keeps water liquid at \SI{300}{\celsius} \\
\bottomrule
\end{tabularx}
\end{table}

% ---------------------------------------------------------------
\section{Mechanical load cases}
% ---------------------------------------------------------------
Each load case prescribes a compressive normal stress on the 16 PEE arcs
(Table~4 of~\cite{Mollaali2026}); PEE $i$ and PEE $i$a carry the same value.
Loads E, B, F and C are Load~A rotated clockwise by \ang{22.5}, \ang{45},
\ang{67.5} and \ang{90} (Table~\ref{tab:loads}).

\begin{table}[htbp]
\centering
\caption{PEE normal stress per load case (MPa, compression positive).}
\label{tab:loads}
\begin{tabular}{@{}lc*{8}{r}@{}}
\toprule
Load & Rotation & PEE 1 & PEE 2 & PEE 3 & PEE 4 & PEE 5 & PEE 6 & PEE 7 & PEE 8 \\
\midrule
A & \ang{0}    & 10.0 & 6.64 & 4.46 & 1.17 & 1.0  & 3.82 & 7.80 & 9.95 \\
E & \ang{22.5} & 9.95 & 10.0 & 6.64 & 4.46 & 1.17 & 1.0  & 3.82 & 7.80 \\
B & \ang{45}   & 7.80 & 9.95 & 10.0 & 6.64 & 4.46 & 1.17 & 1.0  & 3.82 \\
F & \ang{67.5} & 3.82 & 7.80 & 9.95 & 10.0 & 6.64 & 4.46 & 1.17 & 1.0  \\
C & \ang{90}   & 1.0  & 3.82 & 7.80 & 9.95 & 10.0 & 6.64 & 4.46 & 1.17 \\
\bottomrule
\end{tabular}
\end{table}

\begin{itemize}
  \item \textbf{DSS:} each DSS carries the mean of its two neighbouring PEEs,
        $\sigma_{\mathrm{DSS},i}=(\sigma_{\mathrm{PEE},i}+\sigma_{\mathrm{PEE},i+1})/2$,
        with PEE~9 $\equiv$ PEE~1.
  \item \textbf{Out-of-plane:} plane strain; no axial load is applied in 2D.
  \item \textbf{Gravity:} none (horizontal section).
  \item \textbf{Load path:} Scenario~I: applied in the first time step and then
        held constant. Scenario~II: ramped linearly over 0--\SI{1800}{s}, then
        held (Section~\ref{sec:schedule-II}).
\end{itemize}

% ---------------------------------------------------------------
\section{Loading schedule, initial and boundary conditions}
\label{sec:schedule}
% ---------------------------------------------------------------
This section defines the time schedule, initial and boundary conditions of the
two thermal scenarios for the intact sample (THM2). The fractured cases of each
scenario use the same schedule and conditions.

\subsection{Scenario I -- heating}
\label{sec:schedule-I}
THM2 loads the intact sample with Load~A, heats it from the outer boundary for
\SI{4000}{s}, injects cold water at the centre for \SI{50}{s} and then shuts in
until \SI{5000}{s} (Table~\ref{tab:phases}, Fig.~\ref{fig:bc}). M1 is the same
model with only the mechanical load, evaluated at the end of the first step.

\begin{table}[htbp]
\centering
\caption{Scenario I: THM2 phases and boundary conditions.}
\label{tab:phases}
\small
\begin{tabularx}{\textwidth}{@{}llLLl@{}}
\toprule
Phase & Time (s) & Thermal & Hydraulic & Mechanical \\
\midrule
0           & $t=0$     & $T=T_0$ & $p=p_0$ & Load A applied \\
1 Heating   & 0--4000   & Outer boundary: 293 $\rightarrow$ \SI{363}{K} linearly over 0--325\,s, then \SI{363}{K} & $p=p_0$ on outer boundary & Load A held \\
2 Injection & 4000--4050 & As phase 1; centre point $T=\SI{293}{K}$ & Centre point: $Q=\SI{1}{l/min}$ per m $=\SI{1.667e-5}{m^3\,s^{-1}}$ per m (\SI{16.67}{g\,s^{-1}} per m) & Load A held \\
3 Shut-in   & 4050--5000 & As phase 1 & No injection & Load A held \\
\bottomrule
\end{tabularx}
\end{table}

\begin{figure}[htbp]
\centering
\includegraphics[width=\textwidth]{fig_bc.pdf}
\caption{Scenario I boundary conditions: Load~A tractions (a) and the thermal
and hydraulic schedule (b).}
\label{fig:bc}
\end{figure}

\textbf{Initial state:} $T_0=\SI{293}{K}$, $p_0=\SI{1}{MPa}$, $\mathbf u=\mathbf 0$
everywhere. Initial effective stress $\alpha_B p_0$ in the rock (0 in the
sheath), so that total stress is zero before loading. Report strains relative to
this state.

\textbf{Boundaries:} the outer boundary is drained ($p=p_0$) and has the
prescribed temperature in all phases. The reference model treats the sheath
hydraulically and thermally like the rock.

\textbf{Time stepping (reference):} \SI{5}{s} to \SI{400}{s}, \SI{40}{s} to
\SI{4000}{s}, \SI{1}{s} during injection, \SI{2}{s} to \SI{4150}{s}, \SI{10}{s}
to \SI{5000}{s}. Teams may use other steps but must hit all output times exactly.

\subsection{Scenario II -- cold-water injection into hot rock}
\label{sec:schedule-II}
Scenario~II follows the schedule of~\cite{Pan2026} (Tables~8 and~9 there): the
mechanical load is ramped up and held, then water at \SI{15}{\celsius} is
injected at the centre for \SI{3600}{s} into the rock at \SI{300}{\celsius}
(Table~\ref{tab:phases-II}, Fig.~\ref{fig:schedule-II}). The simulation time is
\SI{5600}{s}. There is no rubber sheath; the PEE/DSS loads act on the rock
surface.

\begin{table}[htbp]
\centering
\caption{Scenario II: common \SI{5600}{s} schedule with suggested maximum time
steps~\cite{Pan2026}.}
\label{tab:phases-II}
\small
\begin{tabularx}{\textwidth}{@{}llLLLl@{}}
\toprule
Stage & Time (s) & Mechanical & Injection & Thermal & $\Delta t_\mathrm{max}$ \\
\midrule
Load ramp            & 0--1800    & Linear ramp to target PEE/DSS loads & $Q=0$ & Uniform $T=T_0$ & \SI{10}{s} \\
Hold                 & 1800--2000 & Hold target load & $Q=0$ & Adiabatic equilibration & \SI{10}{s} \\
Cold-water ramp      & 2000--2120 & Hold & Linear $0\rightarrow Q_0$ & \SI{15}{\celsius} inlet & \SI{1}{s} \\
Constant injection   & 2120--5600 & Hold & $Q=Q_0$ & \SI{15}{\celsius} inlet & \SI{5}{s}; $\le$\SI{0.1}{s} at fracture events \\
\bottomrule
\end{tabularx}
\end{table}

\begin{figure}[htbp]
\centering
\includegraphics[width=\textwidth]{fig_schedule_II.pdf}
\caption{Scenario II schedule: mechanical load (fraction of the target PEE/DSS
loads, grey) and injection rate (green).}
\label{fig:schedule-II}
\end{figure}

\textbf{Injection rate:} $Q_0=\SI{0.1}{ml/min}=\SI{1.667e-9}{m^3\,s^{-1}}$
for \SI{3600}{s}, i.e.\ a nominal cumulative volume of \SI{6.0}{ml} (\SI{5.9}{ml}
with the \SI{120}{s} ramp). The 3D rate is converted to the 2D source strength
with the specimen thickness $B=\SI{0.2}{m}$:
\begin{equation}
  Q_\mathrm{2D} = \frac{Q_0}{B} = \SI{8.333e-9}{m^3\,s^{-1}} \text{ per m}
  \quad (\text{mass rate } \SI{8.333e-6}{kg\,s^{-1}} \text{ per m}).
\end{equation}
$Q_0$ must not be applied directly as a 2D line rate.

\begin{table}[htbp]
\centering
\caption{Scenario II: initial and boundary conditions~\cite{Pan2026}.}
\label{tab:icbc-II}
\small
\begin{tabularx}{\textwidth}{@{}lLLL@{}}
\toprule
Field / location & Initial condition & Boundary condition & Comment \\
\midrule
Temperature, whole domain & Uniform rock and pore water at $T_0=\SI{300}{\celsius}$ & Adiabatic outer boundary, $-\mathbf n\cdot(\lambda_\mathrm{eff}\nabla T)=0$ & $T_\mathrm{ref}=T_0$ \\
Pressure, whole domain    & $p=0$ (gauge) & $p=0$ on the outer boundary & Report gauge pressure; $p_\mathrm{abs}=p+\SI{8.6}{MPa}$ \\
Saturation                & $S_w=1$ & Single-phase liquid water & No gas phase \\
Injection point / boundary & $Q=0$ & $Q(t)$ from Table~\ref{tab:phases-II}; switch at the pressure limit & $T=\SI{15}{\celsius}$ only during inflow ($\mathbf q\cdot\mathbf n<0$) \\
Mechanical, outer boundary & Zero load, zero displacement & PEE/DSS loads ramped 0--\SI{1800}{s}, then held & Rigid-body points as in Table~\ref{tab:geometry}; no initial stress \\
Fracture / borehole       & Aperture, permeability and phase-field state as in the HM benchmark & Coupled hydraulic, mechanical and thermal interface & Do not refit \\
\bottomrule
\end{tabularx}
\end{table}

\textbf{Isothermal HM control:} each THM run is repeated with uniform
$T=\SI{15}{\celsius}$ in the whole domain and the injected water, with the same
mesh, time steps and flow schedule (Section~\ref{sec:scenarios}).

\textbf{Fracture cases:} for THM4 and THM5 the breakdown pressure is an output;
a pressure safety limit must not replace the fracture-initiation criterion.

\textbf{Inlet temperature (to be clarified):} a fixed temperature of
\SI{15}{\celsius} at a point source acts as a heat sink that can remove much
more heat than the injected water carries. In the UFZ reference model at the
specified rate it removed about 84 times the heat capacity of the injected
water, so the cooling became conduction-controlled. An energy-consistent inlet
condition (advective enthalpy flux of the injected water) gives much less
cooling. Teams must report which inlet condition they use
(Section~\ref{sec:open}).

% ---------------------------------------------------------------
\section{Required outputs}
\label{sec:outputs}
% ---------------------------------------------------------------
Both scenarios share the sign convention, units and model description below.
The output sets differ: Scenario~I uses the CSV files of version~1
(Section~\ref{sec:outputs-I}); Scenario~II uses the four-point histories,
snapshots, profiles and inlet data of~\cite{Pan2026}
(Section~\ref{sec:outputs-II}).

\textbf{Sign convention:} tension positive for stress and strain in the output
files (the PEE loads are given as compression positive).
$\varepsilon_\mathrm{vol}=\varepsilon_{xx}+\varepsilon_{yy}$;
$\varepsilon_{\theta\theta}$ (\texttt{eps\_tt}) is the circumferential strain,
comparable with the fibre-optic measurement. Angles $\theta$ follow
Section~\ref{sec:angles}.

\textbf{Model description:} each team also submits code and version, process
formulation, element type, mesh size and node count, time stepping, inlet
condition (Scenario~II) and any deviation from this specification.

\subsection{Scenario I}
\label{sec:outputs-I}
Each team delivers four CSV files per rock and load case, in SI units,
comma-separated, with one header line (Table~\ref{tab:outputs}). File names
follow \texttt{<team>\_<case>\_<rock>\_load<X>\_<output>.csv}, e.g.\
\texttt{UFZ\_THM2\_greywacke\_loadA\_probe.csv}.

\begin{table}[htbp]
\centering
\caption{Scenario I: required output files.}
\label{tab:outputs}
\small
\begin{tabularx}{\textwidth}{@{}lLLL@{}}
\toprule
Output & Location & Times & Columns \\
\midrule
\texttt{centre}  & Centre point $(0,0)$ & Every time step, at least every 5\,s; every 1\,s in 4000--4060\,s & \texttt{t, T, p} \\
\texttt{probe}   & Probe circle $r=\SI{0.065}{m}$, 360 points at \ang{1} & Output times & \texttt{t, theta\_deg, eps\_xx, eps\_yy, eps\_xy, eps\_vol, eps\_tt, T, p} \\
\texttt{profile} & Radial line $\theta=\ang{0}$, $0\le r\le\SI{0.1}{m}$, 101 points & Output times & \texttt{t, r, T, p, sigma\_rr, sigma\_tt} \\
\texttt{mean}    & Probe circle average & Every time step & \texttt{t, eps\_vol\_mean, T\_mean, p\_mean} \\
\bottomrule
\end{tabularx}
\end{table}

\textbf{Output times:} 5, 325, 2000, 4000, 4001, 4005, 4010, 4025, 4050, 4060,
4100, 4200, 4500, 5000\,s. For M1 only the state after loading is required
($t=\SI{5}{s}$ or the end of the mechanical step).

\subsection{Scenario II}
\label{sec:outputs-II}
The outputs of Scenario~II follow Tables~10--12 of~\cite{Pan2026}. They are
required for every THM run and for its isothermal HM control.

\textbf{Monitoring points.} Temperature, pore pressure and volumetric strain are
recorded at four points on the probe circle $r_p=\SI{0.065}{m}$
(Table~\ref{tab:points-II}, Fig.~\ref{fig:angles}a). Temperature and pressure
histories must use the same probe circle, coordinate system and interpolation
rule as the volumetric strain.

\begin{table}[htbp]
\centering
\caption{Scenario II: monitoring points on the probe circle~\cite{Pan2026}.}
\label{tab:points-II}
\small
\begin{tabular}{@{}llll@{}}
\toprule
Point & Direction $\theta$ & Coordinates (m) & Required histories \\
\midrule
P0   & \ang{0}   & $(+0.065,\ 0)$ & \multirow{4}{*}{$T(t)$, $p(t)$, $\varepsilon_\mathrm{vol}(t)$} \\
P90  & \ang{90}  & $(0,\ +0.065)$ & \\
P180 & \ang{180} & $(-0.065,\ 0)$ & \\
P270 & \ang{270} & $(0,\ -0.065)$ & \\
\bottomrule
\end{tabular}
\end{table}

\textbf{Required results.} Table~\ref{tab:results-II} lists the results: the
outputs inherited from the HM benchmark, the new thermal outputs and the
injection data.

\begin{table}[htbp]
\centering
\caption{Scenario II: required results~\cite{Pan2026}.}
\label{tab:results-II}
\small
\begin{tabularx}{\textwidth}{@{}lLLl@{}}
\toprule
Category & Variable & Location / form & Unit \\
\midrule
HM (inherited) & Displacement magnitude $|\mathbf u|$, stress trace $\sigma_\mathrm{trace}$, strain trace $\varepsilon_\mathrm{vol}$ & Full field; probe circle & m, MPa, -- \\
HM (inherited) & Pore pressure $p$, Darcy flux $\mathbf q$ & Full field; inlet and outer boundary & MPa, m\,s$^{-1}$ \\
THM (new)      & Temperature $T$ & Full field; P0, P90, P180, P270 & \si{\celsius} \\
THM (new)      & Pore pressure $p(t)$ & P0, P90, P180, P270 & MPa (gauge) \\
THM (new)      & Volumetric strain $\varepsilon_\mathrm{vol}(t)$ & P0, P90, P180, P270 & -- or \si{\micro\epsilon} \\
Inlet          & Inlet pressure $p_\mathrm{in}(t)$, rate $Q(t)$, cumulative volume & Injection point or boundary & MPa, ml/min, ml \\
Fracture (THM4, THM5) & Breakdown pressure and time & Inlet & MPa, s \\
\bottomrule
\end{tabularx}
\end{table}

\textbf{Sampling.} Table~\ref{tab:sampling-II} gives the output times.
Output is refined automatically around fracture events.

\begin{table}[htbp]
\centering
\caption{Scenario II: output sampling~\cite{Pan2026}.}
\label{tab:sampling-II}
\small
\begin{tabularx}{\textwidth}{@{}lL@{}}
\toprule
Output & Sampling \\
\midrule
Four-point histories & From 2000 to \SI{5600}{s}, every time step (at least every \SI{5}{s}, every \SI{1}{s} during the injection ramp); refined around fracture events. Recommended: also 0--\SI{2000}{s} to document the load ramp \\
Inlet data           & As four-point histories \\
Full-field snapshots & $t=2000$, 2060, 3800, \SI{5600}{s}, plus one frame immediately before and after breakdown (THM4, THM5) \\
Circumferential profiles & $T(\theta)$, $p(\theta)$, $\varepsilon_\mathrm{vol}(\theta)$ on $r=r_p$ for $\theta=\ang{0}$--\ang{360} at \ang{1} spacing, at the snapshot times \\
\bottomrule
\end{tabularx}
\end{table}

\textbf{File format.} \cite{Pan2026} prescribes the content but no file format.
To keep both scenarios comparable, Scenario~II uses CSV files with one header
line, named \texttt{<team>\_<case ID>\_<output>.csv} with the case ID of
Section~\ref{sec:scenarios}, e.g.\ \texttt{UFZ\_THM2-GW-A-T300\_points.csv}
(Table~\ref{tab:files-II}). Full-field snapshots are delivered as VTU files or
in another format agreed with the task leads. Temperature in
\si{\celsius}, pressure in MPa (gauge), strain as a fraction, flow rate in
ml/min, volume in ml, all other quantities in SI units.

\begin{table}[htbp]
\centering
\caption{Scenario II: proposed output files per case (THM run and HM control).}
\label{tab:files-II}
\small
\begin{tabularx}{\textwidth}{@{}lLL@{}}
\toprule
File & Content & Columns \\
\midrule
\texttt{points}   & Four-point histories & \texttt{t, T\_P0, p\_P0, eps\_vol\_P0, \ldots, T\_P270, p\_P270, eps\_vol\_P270} \\
\texttt{inlet}    & Injection data & \texttt{t, p\_in, T\_in, Q\_ml\_min, V\_cum\_ml} \\
\texttt{profiles} & Circumferential profiles at the snapshot times & \texttt{t, theta\_deg, T, p, eps\_vol} \\
\texttt{field\_t<time>.vtu} & Full-field snapshot & $|\mathbf u|$, $\sigma_\mathrm{trace}$, $\varepsilon_\mathrm{vol}$, $p$, $\mathbf q$, $T$ \\
\texttt{events}   & THM4, THM5 only & \texttt{t\_breakdown, p\_breakdown} \\
\bottomrule
\end{tabularx}
\end{table}

% ---------------------------------------------------------------
\section{Verification checks}
\label{sec:verification}
% ---------------------------------------------------------------
Each team should pass the checks in Table~\ref{tab:checks} before submitting;
reference values are from the OGS~6.5.9 model of Scenario~I. Checks and
reference values for Scenario~II are still to be added (Section~\ref{sec:open}).

\begin{table}[htbp]
\centering
\caption{Verification checks and reference values ($\varepsilon_\mathrm{vol}$ in $10^{-2}$\,\%).}
\label{tab:checks}
\small
\begin{tabularx}{\textwidth}{@{}lLLLl@{}}
\toprule
Check & Quantity & Reference (OGS) & Independent check & Target \\
\midrule
T: ramp only & $T(r,t)$ & -- & Duhamel solution, disc $R=\SI{0.1}{m}$ & $\le$ 0.2\,K \\
Heating   & Centre $T$ at 4000\,s & 356.7\,K & T model & $\le$ 0.5\,K \\
Injection & $p-p_0$ at $r=\SI{0.065}{m}$ & 0.0812\,MPa & Thiem: 0.0811\,MPa & $\le$ 1\,\% \\
Injection & $p-p_0$ at centre node & 1.12\,MPa & Mesh-dependent & Report only \\
Cooling   & Centre $T$ at 4100 / 5000\,s & 309.5 / 353.6\,K & -- & $\le$ 1\,K \\
M1        & Mean $\varepsilon_\mathrm{vol}$ on probe circle & $-2.46$ (Greywacke), $-0.95$ (Gneiss) & -- & $\le$ 2\,\% \\
Heating   & Mean $\varepsilon_\mathrm{vol}$ at 4000\,s & 11.15 (Greywacke), 12.03 (Gneiss) & $2(1+\nu)\alpha_T\Delta T+\varepsilon_\mathrm{mech}$, $\Delta T=\SI{67.0}{K}$ & $\le$ 2\,\% \\
Injection & Min.\ mean $\varepsilon_\mathrm{vol}$ after injection & 10.73 at 4460\,s (Greywacke), 11.62 at 4450\,s (Gneiss) & -- & Report \\
Recovery  & Mean $\varepsilon_\mathrm{vol}$ at 5000\,s & 10.94 (Greywacke), 11.82 (Gneiss) & -- & $\le$ 2\,\% \\
\bottomrule
\end{tabularx}
\end{table}

The strain on the probe circle varies with angle, with maxima near \ang{100}
and \ang{280} under Load~A. Temperature and pressure are the same for both
rocks, because permeability and bulk thermal properties are equal.

% ---------------------------------------------------------------
\section{Modelling choices and open items}
\label{sec:open}
% ---------------------------------------------------------------
\subsection{Modelling choices to be reported}
Teams must state how they handle the following, since the specification leaves
them open or the reference model uses a simplification:
\begin{itemize}
  \item \textbf{Coupling scheme:} monolithic or staggered; convergence tolerances.
  \item \textbf{Mass vs.\ volume balance:} with pressure-dependent water
        density, injection is a mass rate: \SI{16.67}{g\,s^{-1}} per m
        (Scenario~I), \SI{8.333e-6}{kg\,s^{-1}} per m (Scenario~II).
  \item \textbf{Point injection:} source at a node vs.\ resolved borehole;
        pressure at the injection point is mesh-dependent.
  \item \textbf{Thermal mixing rule} (Scenario~I) used to reach the prescribed
        bulk $\lambda$ and $\rho c$; Scenario~II prescribes
        Eq.~\eqref{eq:mixing}.
  \item \textbf{Inlet temperature condition} (Scenario~II), see
        \subsection{Open items}
The following items are open for the task group. They are grouped by topic;
items marked with a section number are discussed there.

\paragraph{Geometry and cases}
\begin{itemize}[label=\open]
  \item Measured PEE and DSS arc widths (currently \ang{18}/\ang{4.5} estimates,
        Section~\ref{sec:geometry}).
  \item THM4 (Scenario~II): orientation of the initial crack, \ang{45}
        in~\cite{Mollaali2026} (HM3c) vs.\ \ang{-45} in~\cite{Pan2026}.
  \item Case naming: the HM benchmarks were named differently at WS5
        (HM3a/b) and WS6 (M2a/b, HM2a/b), and~\cite{Pan2026} refers to HM3d
        as ``HM5''. Use the mapping of Section~\ref{sec:fracture}
        (THM2--HM1, THM3a--HM2a, THM3b--HM2b, THM4--HM3c, THM5--HM3d)
        consistently in all team contributions.
  \item Scenario~I fractured cases (HM3, HM4, THM3, THM4 of
        Fig.~\ref{fig:concept}): geometry of the fracture network, fracture
        properties and injection into the fracture.
\end{itemize}

\paragraph{Material and fracture parameters}
\begin{itemize}[label=\open]
  \item Gneiss tensile strength: \cite{Mollaali2026} gives 16.8\,MPa only; the
        source and meaning of the second value 6.2\,MPa in~\cite{Pan2026} are
        unclear.
  \item Gneiss thermal expansion: \num{8e-6} (Scenario~I) vs.\
        \num{7e-6}\,K$^{-1}$ (Scenario~II).
  \item Thermal properties of Scenario~I: keep the bulk values
        ($\lambda=2.8$, $\rho c=2700\times790$) or adopt the solid values and
        the mixing rule of Scenario~II (Eq.~\eqref{eq:mixing}).
  \item Thermal properties of Scenario~II are first-round proposals
        of~\cite{Pan2026}; $\lambda_s$, $c_s$ and $\alpha_T$ are to be replaced
        by laboratory measurements when available, without changing case IDs
        or outputs.
  \item Water conductivity at \SI{300}{\celsius}: 0.6 (\cite{Pan2026}) vs.\
        0.553\,W\,m$^{-1}$\,K$^{-1}$ (IAPWS-IF97).
  \item Water compressibility and thermal expansion: \cite{Pan2026} allows
        $\beta_p=\beta_f=0$ up to \SI{100}{\celsius}, but the CAS team argued at
        WS6 against setting them to zero; Scenario~I neglects the thermal
        expansion of water. Agree on one treatment for both scenarios.
  \item Sheath permeability (Scenario~I): \cite{Mollaali2026} gives
        \SI{1e-19}{m^2}; the reference model uses the rock permeability to keep
        a drained outer boundary.
  \item Fractured cases (Section~\ref{sec:fracture}): the factor $\xi$, the
        fracture porosity $\phi_f$ and the choice of AT1 or AT2 are not given
        in~\cite{Mollaali2026}; its Table~2 lists the LIE shear stiffness for
        intact samples as 100/200 instead of 50/100\,GPa\,m$^{-1}$.
\end{itemize}

\paragraph{Schedule, initial and boundary conditions}
\begin{itemize}[label=\open]
  \item Scenario~II inlet temperature: fixed \SI{15}{\celsius} at the inlet
        vs.\ advective enthalpy flux (Section~\ref{sec:schedule-II}).
  \item Scenario~II injection pressure limit: value not given.
  \item Scenario~II injection ramp: this specification and~\cite{Pan2026} use
        2000--\SI{2120}{s}; the CAS results shown at WS6 used 2000--\SI{2020}{s}.
  \item Scenario~II pressure level: the absolute reference \SI{8.6}{MPa} is only
        just above the saturation pressure of water at \SI{300}{\celsius}
        (\SI{8.59}{MPa}); any pressure drop below the initial gauge pressure
        (e.g.\ by cooling contraction of the water at the start of injection)
        violates the single-phase assumption. The absolute reference pressure
        of the \SI{15}{\celsius} HM control is not stated.
  \item Scenario~II fractured cases: hydraulic boundary conditions at the
        fracture ends ($p=\SI{3.45}{MPa}$ in HM2a/b), and whether the common
        rate $Q_\mathrm{2D}$ is sufficient to propagate fractures in THM4 and
        THM5 (HM3c/d used 6 and 45 times higher rates).
  \item Scenario~I: injection rate per metre in 2D vs.\ the \SI{0.2}{m} sample
        height (Scenario~II divides by $B=\SI{0.2}{m}$).
\end{itemize}

\paragraph{Outputs and verification}
\begin{itemize}[label=\open]
  \item Agree the output formats and file naming of Section~\ref{sec:outputs},
        including the Scenario~II files (Table~\ref{tab:files-II}) and the
        format of full-field snapshots.
  \item Angle convention of team results: profiles plotted in the laboratory
        convention appear mirrored and shifted by \ang{22.5}
        (Section~\ref{sec:angles}); the CAS profiles at WS5 and WS6 showed such
        a shift. All teams convert to $\theta$ before submission.
  \item Scenario~II verification checks and reference values
        (Section~\ref{sec:verification}).
  \item The OGS web benchmark curve for HM2a (Load~C) is about
        \SI{26}{\micro\epsilon} less compressive than the reference
        of~\cite{Mollaali2026}; check before it is used as a reference.
  \item Comparison data from GREAT cell heating tests and fibre-optic strain.
\end{itemize}

\paragraph{Document}
\begin{itemize}[label=\open]
  \item Governing equations and the list of participating teams and methods
        (Section~3 and Table~13 of~\cite{Pan2026}) are not yet included.
  \item Corrections to~\cite{Pan2026}: Eqs.~2, 6, 10 and~11 (density and flux
        terms); method labels of UFZ and KIGAM in Table~13 (VPF--FEM,
        GBM--DEM); CAS method named hCA/xFEM in the introduction but
        CASRock.PF in Table~13; LBNL in the nomenclature but not in
        Table~13; cross-references to Table~1 and Section~2.x; labels 2A/3A
        in Fig.~3.
  \item Bibliographic details of~\cite{Mollaali2026} (volume, article number)
        to be completed.
\end{itemize}

 tests and fibre-optic strain.
\end{itemize}

% ---------------------------------------------------------------
\nocite{McDermott2018,Kolditz2025,Mollaali2026}
\bibliographystyle{unsrtnat}
\bibliography{../references}

\end{document}
